\chapter{核最小均方的在线函数更新}
\label{chap:klms}
\begin{quote}\small
\textit{本章摘要。}核最小均方算法在每次收到标签后沿瞬时平方误差的函数梯度更新模型。本章推导核展开、步长与误差之间的关系，再解释字典为何增长，以及如何用新颖性和投影限制字典大小。
\end{quote}

核岭回归对一批训练数据求解线性方程。设备读数连续到达时，反复分解不断增大的核矩阵会增加计算开销。核最小均方算法（Kernel Least Mean Squares，KLMS）在标签到达后只用当前样本更新函数。它适用于平方损失和逐样本更新的场景。本章依据 Theodoridis 第 11.12.1 节\cite{theodoridis2015ml}，将函数空间中的更新写成可执行的核展开。KLMS 是在线随机梯度方法。它不会在每一步重新求解批量核岭回归。

\section{在线学习首先要固定信息顺序}
设时间$t$到达输入$\bm x_t$。先用已有模型$f_{t-1}$预测$\hat y_t=f_{t-1}(\bm x_t)$，随后在$y_t$可用时计算误差并更新。若标签代表一小时后的温度，更新也必须等到相应标签确认。用于评价的$\hat y_t$应保留为更新前预测，不能用已经看过$y_t$的新函数重新预测同一样本来报告在线误差。

在平稳分布下，可以考虑总体风险
\[
 R(f)=\frac12\E[(Y-f(X))^2].
\]
其中$(X,Y)$服从数据分布，期望需要对该分布计算。这个分布通常未知，所以不能直接得到总体梯度。KLMS用当前样本的瞬时损失$\ell_t(f)=(y_t-f(\bm x_t))^2/2$构造随机梯度。在独立同分布抽样和适当可积条件下，这种梯度在期望意义上对应总体梯度。时间序列存在相关性或分布漂移时，必须重新检查这一解释。

\section{函数梯度怎样变成一个核截面}
由再生性质$f(\bm x_t)=\langle f,k(\cdot,\bm x_t)\rangle$。沿任意方向$h\in\mathcal H$作微小变化，得到
\begin{equation}
 \left.\frac{d}{ds}\ell_t(f+sh)\right|_{s=0}
 =-(y_t-f(\bm x_t))h(\bm x_t)
 =\langle-e_t k(\cdot,\bm x_t),h\rangle,
\end{equation}
其中$e_t=y_t-f(\bm x_t)$。这说明函数梯度为$-e_t k(\cdot,\bm x_t)$。用步长$\mu_t>0$沿负梯度更新：
\begin{equation}
 e_t=y_t-f_{t-1}(\bm x_t),\qquad
 f_t=f_{t-1}+\mu_t e_t k(\cdot,\bm x_t).
 \label{eq:klms-update}
\end{equation}
误差为正时提高附近输入的预测，误差为负时降低。影响如何随输入变化由核值决定，有效核并不要求所有核值都非负。

从$f_0=0$开始展开递推，得到
\begin{equation}
 f_t(\bm x)=\sum_{s=1}^t\mu_s e_s k(\bm x_s,\bm x),\qquad
 \hat y_t=\sum_{s=1}^{t-1}\mu_s e_s k(\bm x_s,\bm x_t).
 \label{eq:klms-expansion}
\end{equation}
每次保存输入$\bm x_t$和系数$\mu_t e_t$即可。既往系数包含当时的误差，不需要根据当前模型重新计算。图\ref{fig:klms-online-update}突出预测与更新的顺序。

\begin{figure}[htbp]
\centering
\begin{tikzpicture}[node distance=0.8cm and 0.7cm]
 \node[idi input] (x) {$\bm x_t,f_{t-1}$};
 \node[idi process,right=of x] (p) {$\hat y_t=f_{t-1}(\bm x_t)$};
 \node[idi state,right=of p] (e) {$e_t=y_t-\hat y_t$};
 \node[idi output,below=of e] (u) {$f_t=f_{t-1}+\mu_te_tk(\cdot,\bm x_t)$};
 \node[idi input,above=of e] (y) {标签$y_t$确认};
 \draw[idi arrow] (x)--(p); \draw[idi arrow] (p)--(e);
 \draw[idi arrow] (y)--(e); \draw[idi arrow] (e)--(u);
\end{tikzpicture}
\caption{KLMS一次在线更新。箭头表示计算依赖，当前预测先于标签参与更新。标签延迟时应保留输入与更新前预测。}
\label{fig:klms-online-update}
\end{figure}

\section{步长怎样改变当前残差}
在同一个已观测输入上，更新后的残差为
\begin{equation}
 y_t-f_t(\bm x_t)=\{1-\mu_t k(\bm x_t,\bm x_t)\}e_t.
\end{equation}
若$k(\bm x_t,\bm x_t)>0$，则$0<\mu_t k(\bm x_t,\bm x_t)<2$使当前点的残差绝对值减小。这只是一次局部更新的结论，不保证未来样本误差减小，也不替代随机迭代的收敛条件。

可使用归一化更新
\begin{equation}
 f_t=f_{t-1}+\frac{\eta e_t}{\delta+k(\bm x_t,\bm x_t)}k(\cdot,\bm x_t),
 \qquad\delta>0,\quad0<\eta<2.
\end{equation}
此时更新后的当前残差为
\[
 y_t-f_t(\bm x_t)=\left(1-\frac{\eta k(\bm x_t,\bm x_t)}{\delta+k(\bm x_t,\bm x_t)}\right)e_t.
\]
当$k(\bm x_t,\bm x_t)>0$时，$0<\eta<2$使这个乘子的绝对值小于1。分母按当前核截面的平方范数调整步长，$\delta$保证分母为正。若令$\delta$趋近于零，乘子趋近于$1-\eta$。对高斯核，$k(\bm x,\bm x)=1$，未归一化更新的局部条件是$0<\mu_t<2$。在平稳随机优化中，常用满足$\sum_t\mu_t=\infty$且$\sum_t\mu_t^2<\infty$的递减步长，并要求梯度噪声和核截面具有适当界。这些条件不能直接推出任意相关工业序列上的收敛。固定步长有利于跟踪变化，但通常保留稳态波动。

\subsection{两次更新的算例}
取高斯核长度尺度$\ell=1$，固定步长$\mu=1/2$。第一对观测为$(x_1,y_1)=(0,1)$，初始预测为0，所以$e_1=1$，$f_1(x)=0.5\exp(-x^2/2)$。第二对为$(x_2,y_2)=(1,0)$，更新前预测为$0.5\exp(-1/2)\approx0.3033$，误差约为$-0.3033$。因此
\begin{equation}
 f_2(x)=0.5k(0,x)-0.1516k(1,x).
\end{equation}
第二点更新后预测约为0.1516，残差幅度减半。第一点预测也从0.5降为约0.4080，因为两个核截面存在非零内积。修正一个点会影响别的点，这解释了局部下降与总体下降的区别。

\section{字典为什么增长，怎样限制规模}
式\eqref{eq:klms-expansion}每次添加一个核中心，处理$t$个样本后单次预测需要$O(t)$核计算，累计成本为$O(t^2)$，存储中心需要$O(td)$。因此在线并不自动意味着恒定时间与内存。实际使用字典$D=\{\bm u_1,\ldots,\bm u_m\}$，只保留一部分核中心。

一种简单新颖性策略同时检查新输入与字典的最小距离以及预测误差：只有距离超过阈值且$|e_t|$超过阈值时才添加中心。它会舍弃部分更新，需要验证带来的误差。原空间距离还依赖缩放，可以改用核诱导距离
\begin{equation}
 \|k(\cdot,\bm x)-k(\cdot,\bm u)\|_{\mathcal H}^2
 =k(\bm x,\bm x)+k(\bm u,\bm u)-2k(\bm x,\bm u).
\end{equation}
若想判断新核截面是否已能被整个字典表示，可将其投影到字典张成的空间。令$G_{ij}=k(\bm u_i,\bm u_j)$，并令$\bm g=(k(\bm u_i,\bm x_t))_{i=1}^m$。在字典核矩阵$G$可逆时，投影系数和剩余平方长度为
\begin{equation}
 \bm c=G^{-1}\bm g,\qquad
 r_D^2=k(\bm x_t,\bm x_t)-\bm g^{\mathsf T}G^{-1}\bm g\ge0.
 \label{eq:klms-dictionary-residual}
\end{equation}
若$r_D^2$很小，可用$\sum_{i=1}^m c_i k(\cdot,\bm u_i)$近似当前核截面，并把$\mu_t e_t\bm c$加到既有字典系数上。若$r_D^2$较大，则添加新中心。$G$奇异时可用伪逆或稳定分解。若计算采用$G+\rho I$这样的正则近似，所得系数不是精确正交投影，不能把对应数值直接解释为式\eqref{eq:klms-dictionary-residual}中的残差。

该投影近似在一次更新中引入的函数范数误差为$|\mu_te_t|r_D$。由再生性质，在查询点$\bm x$的预测误差至多为$|\mu_te_t|r_D\sqrt{k(\bm x,\bm x)}$。这给出了可检查的局部近似含义，但多步误差仍会累积。

\section{正则化更新与评价边界}
若瞬时目标另加$\lambda\|f\|^2/2$，函数梯度多出$\lambda f$，得到
\begin{equation}
 f_t=(1-\mu_t\lambda)f_{t-1}+\mu_te_tk(\cdot,\bm x_t).
\end{equation}
它缩小既有系数，属于正则化在线梯度更新，可与Theodoridis第11.12.2节的朴素在线正则风险最小化算法（Naive Online Regularized Risk Minimization Algorithm，NORMA）对照。只有系数变小还不会使字典有限，仍需剪枝或预算策略。这是对KLMS目标的修改，不能将其正则化参数直接视为批量KRR的同一个解。

在线评价应按时间先预测后更新，并记录标签延迟、字典规模、单次更新时间和分布漂移后的恢复过程。打乱样本顺序适合部分平稳优化实验，却不能模拟真实设备的时间部署。核尺度、步长和字典阈值应在历史验证数据上确定。测试数据不能被用来逐次调整这些设置。

\section{阅读依据}
Theodoridis第11.12.1节（书内页553--556）提供KLMS、归一化更新与字典思想，第11.12.2节（页556--560）提供正则化在线更新。本章在这些基础上明确局部步长条件与总体收敛的区别，并补充字典投影误差的推导。
